\providecommand{\simulationdir}{coder}
\section{Simulation of the construction}
\label{sec:simulation}

We now test the local correlation calculation by sampling point coordinates
from the construction. We use $M=1$ and $\lambda=6\times10^{-6}$, with the
explicit parameter choices in the proof. These give
$\lambda_C=\lambda_Y=\lambda$, $L=1$,
$K\simeq1.75519\times10^{-8}$ and $c\simeq9.82728\times10^{-10}$.
The small intensity makes direct sampling instructive, but inefficient.
In five million independent realisations restricted to the ball $B(0,0.45)$,
we observed seven points, each in a different realisation. Four came from
the sphere alternative and three from the correction alternative. The
expected total point count is $11.4511$. Thus no pair, triangle or tetrahedron
was observed in these ordinary windows. We keep these data separate from
the importance-sampled experiment below.

\paragraph{Exact simulation in a small ball.}
The separation in the construction means that a ball of diameter less than
$M$ meets at most one retained sphere and receives points from at most one
correction cluster. Its restricted law is therefore determined by the
first measure of the possible intersecting object. For the sphere component,
this measure is $\gamma(R)\,dz\,dR$, restricted to spheres meeting the ball.
If the ball has radius $r$, we dominate its centre domain by $B(0,R+r)$.
The total dominating mass is
\[
 p_* = \frac{8\pi F}{3L^4}
        \int_0^\infty R(R+r)^3 e^{-4R^2/L^2}\,dR.
\]
For $r=0.45$ and our parameters, $p_*\simeq9.04762\times10^{-5}<1$.
With probability $p_*$ we draw $R$ from the normalised integrand and a
centre uniformly from $B(0,R+r)$; otherwise the window is empty. Expanding
$R(R+r)^3$ gives a four-component mixture in which $R^2$ has a gamma
distribution, so no radius tail needs to be truncated. We reject a sphere
that does not meet the ball. For an intersecting sphere we draw the
Poisson count with mean $a(R)/(2R)$ times its intersection area, followed
by independent uniform points on that spherical cap or whole sphere.

The correction component admits the same single-object method, using
centres in $B(0,r+1)$, their intensity $\kappa$, and the actual cluster
mark law. The auxiliary rejection construction described below samples its
marks without evaluating the marginal $\mathcal Mq_3$ numerically.
We flip the fair whole-law coin once for each independent window. This
algorithm produces exact restrictions of the stated process, up to
floating-point implementation; it does not construct several coupled
windows from one whole-space realisation. For a general window one can
instead sample the sphere band, shifts and finitely many relevant candidate
spheres directly. The correction alternative can be sampled from Poisson
parents on the window enlarged by $4M$, testing the $3M$ exclusion condition
for parents within distance $M$ of the window.

\paragraph{The measured correlations.}
Write $g_n=\rho^{(n)}/\lambda^n$. We anchor the first vertex at zero and
let $H_{nj}$ indicate a simplex-size bin, together with the condition that
all pair distances are below 1. At order four we also require the unique
circumsphere radius to be at most 2. Each reported bin has a positive lower
edge. Our estimand is
\begin{equation}
 \overline g_{nj}=
 \frac{\int H_{nj}(y)\rho^{(n)}(0,y)\,dy}
      {\lambda^n\int H_{nj}(y)\,dy}.
 \label{eq:simulation-estimand}
\end{equation}
These restrictions make finite observables; they do not change the process
or the domains of the theorem. The theoretical comparison is the bin average
of $a_n/(\lambda^n\Delta_n)$, not its value at a bin midpoint. Anchored
relative coordinates avoid window boundary loss: by stationarity, the
numerator is the expected number of such ordered tuples whose first point
lies in a unit-volume observation set, with the remaining points unrestricted.
We use the prescribed $\lambda$ throughout, since a spatial intensity
estimate from one nonergodic realisation need not estimate this mixture mean.

\paragraph{Importance-sampled sphere tuples.}
Let $b(R)=a(R)/(2R)$ and $S(R)=\pi\min(1,4R^2)$, the area of the cap of a
sphere through zero lying in $B(0,1)$. Integrating the first point's surface
measure over its possible centres gives
\begin{equation}
 \int H_{nj}r_n
 =\int_0^\infty 4\pi R^2\gamma(R)b(R)^n S(R)^{n-1}
       E\{H_{nj}(0,X_2,\ldots,X_n)\mid R\}\,dR.
 \label{eq:simulation-sphere-sampling}
\end{equation}
Here the centre is $Re$, where $e$ is uniform on the unit sphere, and the
remaining points are independent uniform points in the cap. This identity
comes directly from the conditional Poisson factorial measure, before its
geometric density is evaluated. It integrates out the Poisson count using
$E(N)_n=m^n$. The sampled rows are factorial tuples, rather than ordinary
Poisson realisations; unused background points are not generated.

We sample $R$ with density
\[
 h_2(R)=\frac{4R}{L^2}e^{-2R^2/L^2},\qquad
 h_3(R)=\frac{2}{L\sqrt\pi}e^{-R^2/L^2}.
\]
The corresponding importance weights in
\eqref{eq:simulation-sphere-sampling} are $2A_2S(R)$ and
$4\pi A_3S(R)^2$. Both proposals retain the complete positive radius range.
For four points the generating sphere is the tuple's unique circumsphere,
so the observable's guard permits $0<R\leq2$. With
$Z=1/6+\log4$, we use $h_4(R)=S(R)^3/(\pi^3ZR)$ on this interval.
Its weight is the constant $\pi^4cZ/2$. All these weights follow from the
generative intensity $\gamma$ and surface intensity $b$; the reciprocal
simplex formula is not inserted into the empirical numerator. Distances,
areas and volumes are calculated from the sampled Cartesian coordinates.

\paragraph{Correction-cluster marks.}
We sample actual two- and three-point cluster marks by rejection, avoiding
a numerical approximation to $F_2=q_2-\mathcal Mq_3$. For $M=1$, put
$\varepsilon=1/(2L)$ and $J=2\varepsilon-\varepsilon^2$. Write
$u=de$ and $v=te+r(\cos\phi\,e_1+\sin\phi\,e_2)$ in a tangent basis.
The integrable envelope
\[
 E(0,u,v)=\frac{2A_3}{dr}\min(1,\varepsilon^2/r^2)
      \mathbf1_{\{0<d<1,\ |t|<1,\ 0<r<1\}}
\]
dominates $q_3$, and its mass is $T=16\pi^2A_3J$.
It is sampled by independent draws of $d$ with density $2d$, uniform $e$,
$t$ and $\phi$, and $r$ with density
$\min(1,\varepsilon^2/r^2)/J$. Acceptance with probability $q_3/E$
produces the exact three-point mark law. Each accepted row contributes
$TH_{3j}$ to the third integral; its contribution to the second integral
is $T/3$ times the sum of the three edge-bin indicators. Symmetry and
re-rooting give the latter identity, and the three edge contributions are
kept together when calculating variances.

For a pair, draw $u$ from $q_2/Q_2$ by radial rejection, then draw an auxiliary
$v$ from the cylindrical conditional proposal $h(v\mid u)$ above. Accept
the pair with probability
\[
 1-\frac{q_3(0,u,v)}{q_2(0,u)h(v\mid u)}.
\]
The subtracted quantity is at most $JL/2\leq1/2$. Integrating over $v$ gives
the unnormalised accepted density $F_2/Q_2$; thus an accepted row contributes
$Q_2H_{2j}$. The auxiliary vertex is not an emitted cluster point.
The retained-parent intensity and the cluster size probabilities cancel
the ordered-selection factors exactly as in the proof. We therefore need
only sample marks for this local factorial experiment; we do not replace
the dependent retained parents by independent ones.

For the ordinary-window implementation, one can avoid $Q_3$ even in the
mark probabilities. At a retained centre, use candidate masses
$(\lambda_C-Q_2)/\kappa$, $Q_2/(2\kappa)$, $T/(6\kappa)$ and
$T/(2\kappa)$. These respectively propose a singleton, an accepted pair,
an accepted triple, or an accepted triple proposal that emits only its
centre. Rejections emit nothing. Their sum is at most one because
$T\leq Q_2/2$ and $\lambda_C\leq\kappa$; their accepted masses are exactly
the prescribed singleton, pair and triple mark probabilities.

\paragraph{Results and uncertainty.}
For each order we generated $524\,288$ independent sphere proposals. We also
generated $524\,288$ pair proposals and $524\,288$ triple proposals,
accepting $497\,387$ pair marks and $90\,470$ triple marks. Independent
reference streams used $2\,097\,152$ uniform relative configurations per
order to integrate bin volumes and reciprocal simplex sizes. These
reference samples are numerical integration data, separately labelled from
the process samples. The full factorial integral estimates are
one half the sphere estimate plus one half the relevant correction estimate,
implementing the whole-law mixture without cross-component tuples.

All process contributions in this experiment are bounded, as are the
reference integrands on the reported bins. Ordinary independent-sample
standard errors therefore apply to these weighted observations, despite
the original infinite fourth count moment. We propagate both the process
sampling error and the independent reference integration error, retaining
the covariance of the reference numerator and denominator. Error bars are
$1.96$ estimated Monte Carlo standard errors; they are approximate
pointwise intervals, not finite-sample coverage guarantees.

The ratios of estimated to theoretical bin averages range from $0.9767$
to $1.0043$. Their standard errors range from $0.0019$ to $0.0253$;
every reported deviation is below two standard errors. The correction
contributes approximately $45$--$64\%$ of the pair integrals and
$17$--$51\%$ of the triangle integrals in these bins. Its fourth contribution
is zero by construction. Figure~\ref{fig:simulation-correlations} and
Table~\ref{tab:simulation-results} report the normalised correlations;
the supplement contains all bins, full precision and separate reference
integration errors. This agreement checks the sampled construction and
the local measure calculation in the stated domains. It does not establish
the exact identities, global regularity, or infinite fourth moment, which
are consequences of the proof.

\begin{figure}[htbp]
 \centering
 \includegraphics[width=\linewidth]{\simulationdir/figures/correlations.pdf}
 \caption{Bin-averaged $g_2,g_3,g_4$ from sampled factorial tuples and
 correction-cluster marks, with independent numerical integration of the
 theoretical bin averages. All bins have diameter below 1; order four
 additionally has $R_{\rm circ}\leq2$. The lower panels show the ratio to
 theory. Error bars include sampling and reference integration error and
 equal $1.96$ Monte Carlo standard errors.}
 \label{fig:simulation-correlations}
\end{figure}

\input{\simulationdir/results_table.tex}

\begin{figure}[htbp]
 \centering
 \includegraphics[width=\linewidth]{\simulationdir/figures/sampled_configurations.pdf}
 \caption{Saved point coordinates from three distinct kinds of sample.
 Left: the seven points from five million ordinary independent small-ball
 realisations are overlaid; they belong to seven different realisations.
 Centre: a separate conditional Poisson illustration, given a sphere of
 radius 2 centred at $(0,0,2)$ and at least four points in $B(0,0.45)$.
 Right: one accepted three-point correction-cluster mark, before translation
 by its parent. The latter two illustrations are not ordinary draws from
 the whole-law mixture and are not added to the ordinary-window counts.}
 \label{fig:simulation-configurations}
\end{figure}

The reproducible supplement contains the code, random seeds, every process
tuple and its contribution, all ordinary-window counts and nonempty point
clouds, and the reference sufficient statistics. A feasibility pilot was
kept separate; it led us to end the reported fourth-order bins at volume
$0.075$, since larger volumes were rare under the proposal. The simulation
does not truncate the theorem's domain, claim ergodicity, or compare the
construction with earthquake catalogues.
